% --- Archon physics formalization source begin ---
% archon:physics
% archon:covers PhyXMiniProblems/problem_phyx_mini_0081.lean
% archon:source-report reports/phyx_mini/problem_phyx_mini_0081.source.json

\chapter{Physics problem phyx\_mini\_0081}
\label{ch:PhyXMiniProblems_problem_phyx_mini_0081}

\paragraph{Problem source.}
\#\# Physical scenario

Figure is a graph of intensity versus angular position \$	heta\$ for the diffraction of an x-ray beam by a crystal. The horizontal scale is set by \$	heta\_s = 2.00\textasciicircum{}\textbackslash{}circ\$. The beam consists of two wavelengths, and the spacing between the reflecting planes is \$0.94 \textbackslash{}, \textbackslash{}text\{nm\}\$.

\#\# Question

What is the longer wavelengths in the beam?

\#\# Answer choices

- A: A: 30pm

- B: B: 35pm

- C: C: 38pm

- D: D: 25pm

\#\# Figure caption (auxiliary; use the image as primary evidence)

The image is a graph plotting intensity against the angle θ (in degrees). 

- **Axes**: 
  - The x-axis is labeled "θ (degrees)" and spans from 0 to an unspecified maximum value, marked as \textbackslash{}( \textbackslash{}theta\_s \textbackslash{}).
  - The y-axis is labeled "Intensity."

- **Plot**: 
  - The graph features a green line showing the variation of intensity with respect to the angle.
  - There are multiple peaks, with the first occurring at a lower angle, featuring a sharp and prominent peak.
  - Subsequent peaks are present, with varying heights and widths.

- **Grid**: The graph has a grid background that helps in visualizing the data points more clearly.

The overall trend shows intensity having several distinct peak points across varying angles.

\#\# Dataset metadata

Wave Optics; Physical Model Grounding Reasoning, Multi-Formula Reasoning

\paragraph{Recorded answer/context.}
C: C: 38pm

\paragraph{Figure/image path.}
/root/proposal\_for\_physic/science-mango/phyx\_mini\_run/phyx\_data/test\_image/81.png

\paragraph{Formalization target.}
create a compiling Lean file with sorry bodies at `PhyXMiniProblems/problem\_phyx\_mini\_0081.lean`. The Lean declarations must preserve the physical quantities, dimensions or dimensional roles, figure labels, governing-law hypotheses, and final relation expressed by this problem.
Use Mathlib/Physlib names found through LeanExplore where available. If a physics API is missing, introduce faithful local abstractions rather than scalar placeholder aliases.

\begin{theorem}[Physics formalization target]
\label{thm:physics:phyx_mini_0081:target}
\lean{PhyXMiniProblems.ProblemPhyXMini0081.problem_phyx_mini_0081}
\uses{def:physics:phyx-mini-0081:phyxminiproblems-problemphyxmini0081-twowavelengthxraybeam, def:physics:phyx-mini-0081:phyxminiproblems-problemphyxmini0081-reflectingcrystal, def:physics:phyx-mini-0081:phyxminiproblems-problemphyxmini0081-angulardiffractionfigure, def:physics:phyx-mini-0081:phyxminiproblems-problemphyxmini0081-matchesproblemandfigurereadouts, def:physics:phyx-mini-0081:phyxminiproblems-problemphyxmini0081-hasphysicaltwolineparameters, def:physics:phyx-mini-0081:phyxminiproblems-problemphyxmini0081-satisfiestwolinebragglaws, def:physics:phyx-mini-0081:phyxminiproblems-problemphyxmini0081-answerwavelengthpicometers, def:physics:phyx-mini-0081:phyxminiproblems-problemphyxmini0081-isuniqueclosestdisplayedanswer}
The assigned autoformalize agent should translate this physics problem into Lean declarations in the covered file, with theorem and lemma proof bodies written as `by sorry`.
\end{theorem}
\begin{proof}
This is an autoformalization task, not a proof task. Produce statements that can later be proved by the physics prover without weakening the physical model.
\end{proof}

% --- Archon declaration topology begin ---
\section{Declaration topology}
\begin{definition}[Length Magnitude]
  \label{def:physics:phyx-mini-0081:phyxminiproblems-problemphyxmini0081-lengthmagnitude}
  \lean{PhyXMiniProblems.ProblemPhyXMini0081.LengthMagnitude}
  A nonnegative physical quantity carrying the dimension of length
\end{definition}
\begin{proof}
  This is the declared model definition of Length Magnitude.
\end{proof}
\begin{definition}[length Value In]
  \label{def:physics:phyx-mini-0081:phyxminiproblems-problemphyxmini0081-lengthvaluein}
  \lean{PhyXMiniProblems.ProblemPhyXMini0081.lengthValueIn}
  \uses{def:physics:phyx-mini-0081:phyxminiproblems-problemphyxmini0081-lengthmagnitude}
  Read a physical length as a real scalar in the specified length unit
\end{definition}
\begin{proof}
  This is the declared model definition of length Value In.
\end{proof}
\begin{definition}[nanometers Value]
  \label{def:physics:phyx-mini-0081:phyxminiproblems-problemphyxmini0081-nanometersvalue}
  \lean{PhyXMiniProblems.ProblemPhyXMini0081.nanometersValue}
  \uses{def:physics:phyx-mini-0081:phyxminiproblems-problemphyxmini0081-lengthmagnitude, def:physics:phyx-mini-0081:phyxminiproblems-problemphyxmini0081-lengthvaluein}
  The nanometer readout used for the crystal spacing and Bragg equation
\end{definition}
\begin{proof}
  This is the declared model definition of nanometers Value.
\end{proof}
\begin{definition}[picometers Value]
  \label{def:physics:phyx-mini-0081:phyxminiproblems-problemphyxmini0081-picometersvalue}
  \lean{PhyXMiniProblems.ProblemPhyXMini0081.picometersValue}
  \uses{def:physics:phyx-mini-0081:phyxminiproblems-problemphyxmini0081-lengthmagnitude, def:physics:phyx-mini-0081:phyxminiproblems-problemphyxmini0081-lengthvaluein}
  The picometer readout used by the displayed answer choices
\end{definition}
\begin{proof}
  This is the declared model definition of picometers Value.
\end{proof}
\begin{definition}[degrees To Radians]
  \label{def:physics:phyx-mini-0081:phyxminiproblems-problemphyxmini0081-degreestoradians}
  \lean{PhyXMiniProblems.ProblemPhyXMini0081.degreesToRadians}
  Convert a scalar angle readout in degrees to radians
\end{definition}
\begin{proof}
  This is the declared model definition of degrees To Radians.
\end{proof}
\begin{definition}[Spectral Line]
  \label{def:physics:phyx-mini-0081:phyxminiproblems-problemphyxmini0081-spectralline}
  \lean{PhyXMiniProblems.ProblemPhyXMini0081.SpectralLine}
  The two spectral components of the incident x-ray beam
\end{definition}
\begin{proof}
  This is the declared model definition of Spectral Line.
\end{proof}
\begin{definition}[Diffraction Peak]
  \label{def:physics:phyx-mini-0081:phyxminiproblems-problemphyxmini0081-diffractionpeak}
  \lean{PhyXMiniProblems.ProblemPhyXMini0081.DiffractionPeak}
  The four diffraction peaks, ordered from left to right in the figure
\end{definition}
\begin{proof}
  This is the declared model definition of Diffraction Peak.
\end{proof}
\begin{definition}[Figure Axis Label]
  \label{def:physics:phyx-mini-0081:phyxminiproblems-problemphyxmini0081-figureaxislabel}
  \lean{PhyXMiniProblems.ProblemPhyXMini0081.FigureAxisLabel}
  Labels printed on the two axes of the primary figure
\end{definition}
\begin{proof}
  This is the declared model definition of Figure Axis Label.
\end{proof}
\begin{definition}[Two Wavelength XRay Beam]
  \label{def:physics:phyx-mini-0081:phyxminiproblems-problemphyxmini0081-twowavelengthxraybeam}
  \lean{PhyXMiniProblems.ProblemPhyXMini0081.TwoWavelengthXRayBeam}
  \uses{def:physics:phyx-mini-0081:phyxminiproblems-problemphyxmini0081-lengthmagnitude, def:physics:phyx-mini-0081:phyxminiproblems-problemphyxmini0081-spectralline}
  An x-ray beam consisting of two physical wavelength components
\end{definition}
\begin{proof}
  This is the declared model definition of Two Wavelength XRay Beam.
\end{proof}
\begin{definition}[Reflecting Crystal]
  \label{def:physics:phyx-mini-0081:phyxminiproblems-problemphyxmini0081-reflectingcrystal}
  \lean{PhyXMiniProblems.ProblemPhyXMini0081.ReflectingCrystal}
  \uses{def:physics:phyx-mini-0081:phyxminiproblems-problemphyxmini0081-lengthmagnitude}
  The family of parallel reflecting planes in the crystal
\end{definition}
\begin{proof}
  This is the declared model definition of Reflecting Crystal.
\end{proof}
\begin{definition}[Angular Diffraction Figure]
  \label{def:physics:phyx-mini-0081:phyxminiproblems-problemphyxmini0081-angulardiffractionfigure}
  \lean{PhyXMiniProblems.ProblemPhyXMini0081.AngularDiffractionFigure}
  \uses{def:physics:phyx-mini-0081:phyxminiproblems-problemphyxmini0081-diffractionpeak, def:physics:phyx-mini-0081:phyxminiproblems-problemphyxmini0081-figureaxislabel}
  This structure records the Angular Diffraction Figure component of the problem-specific physical model
\end{definition}
\begin{proof}
  This is the declared model definition of Angular Diffraction Figure.
\end{proof}
\begin{definition}[spectral Line At Peak]
  \label{def:physics:phyx-mini-0081:phyxminiproblems-problemphyxmini0081-spectrallineatpeak}
  \lean{PhyXMiniProblems.ProblemPhyXMini0081.spectralLineAtPeak}
  \uses{def:physics:phyx-mini-0081:phyxminiproblems-problemphyxmini0081-spectralline, def:physics:phyx-mini-0081:phyxminiproblems-problemphyxmini0081-diffractionpeak}
  This def records the spectral Line At Peak component of the problem-specific physical model
\end{definition}
\begin{proof}
  This is the declared model definition of spectral Line At Peak.
\end{proof}
\begin{definition}[diffraction Order At Peak]
  \label{def:physics:phyx-mini-0081:phyxminiproblems-problemphyxmini0081-diffractionorderatpeak}
  \lean{PhyXMiniProblems.ProblemPhyXMini0081.diffractionOrderAtPeak}
  \uses{def:physics:phyx-mini-0081:phyxminiproblems-problemphyxmini0081-diffractionpeak}
  The corresponding first- or second-order Bragg reflection
\end{definition}
\begin{proof}
  This is the declared model definition of diffraction Order At Peak.
\end{proof}
\begin{definition}[Matches Problem And Figure Readouts]
  \label{def:physics:phyx-mini-0081:phyxminiproblems-problemphyxmini0081-matchesproblemandfigurereadouts}
  \lean{PhyXMiniProblems.ProblemPhyXMini0081.MatchesProblemAndFigureReadouts}
  \uses{def:physics:phyx-mini-0081:phyxminiproblems-problemphyxmini0081-nanometersvalue, def:physics:phyx-mini-0081:phyxminiproblems-problemphyxmini0081-reflectingcrystal, def:physics:phyx-mini-0081:phyxminiproblems-problemphyxmini0081-angulardiffractionfigure}
  This def records the Matches Problem And Figure Readouts component of the problem-specific physical model
\end{definition}
\begin{proof}
  This is the declared model definition of Matches Problem And Figure Readouts.
\end{proof}
\begin{definition}[Has Physical Two Line Parameters]
  \label{def:physics:phyx-mini-0081:phyxminiproblems-problemphyxmini0081-hasphysicaltwolineparameters}
  \lean{PhyXMiniProblems.ProblemPhyXMini0081.HasPhysicalTwoLineParameters}
  \uses{def:physics:phyx-mini-0081:phyxminiproblems-problemphyxmini0081-nanometersvalue, def:physics:phyx-mini-0081:phyxminiproblems-problemphyxmini0081-picometersvalue, def:physics:phyx-mini-0081:phyxminiproblems-problemphyxmini0081-twowavelengthxraybeam, def:physics:phyx-mini-0081:phyxminiproblems-problemphyxmini0081-reflectingcrystal}
  This def records the Has Physical Two Line Parameters component of the problem-specific physical model
\end{definition}
\begin{proof}
  This is the declared model definition of Has Physical Two Line Parameters.
\end{proof}
\begin{definition}[Satisfies Two Line Bragg Laws]
  \label{def:physics:phyx-mini-0081:phyxminiproblems-problemphyxmini0081-satisfiestwolinebragglaws}
  \lean{PhyXMiniProblems.ProblemPhyXMini0081.SatisfiesTwoLineBraggLaws}
  \uses{def:physics:phyx-mini-0081:phyxminiproblems-problemphyxmini0081-nanometersvalue, def:physics:phyx-mini-0081:phyxminiproblems-problemphyxmini0081-degreestoradians, def:physics:phyx-mini-0081:phyxminiproblems-problemphyxmini0081-twowavelengthxraybeam, def:physics:phyx-mini-0081:phyxminiproblems-problemphyxmini0081-reflectingcrystal, def:physics:phyx-mini-0081:phyxminiproblems-problemphyxmini0081-angulardiffractionfigure, def:physics:phyx-mini-0081:phyxminiproblems-problemphyxmini0081-spectrallineatpeak, def:physics:phyx-mini-0081:phyxminiproblems-problemphyxmini0081-diffractionorderatpeak}
  This structure records the Satisfies Two Line Bragg Laws component of the problem-specific physical model
\end{definition}
\begin{proof}
  This is the declared model definition of Satisfies Two Line Bragg Laws.
\end{proof}
\begin{definition}[Answer Choice]
  \label{def:physics:phyx-mini-0081:phyxminiproblems-problemphyxmini0081-answerchoice}
  \lean{PhyXMiniProblems.ProblemPhyXMini0081.AnswerChoice}
  Labels of the four answers displayed with the problem
\end{definition}
\begin{proof}
  This is the declared model definition of Answer Choice.
\end{proof}
\begin{definition}[answer Wavelength Picometers]
  \label{def:physics:phyx-mini-0081:phyxminiproblems-problemphyxmini0081-answerwavelengthpicometers}
  \lean{PhyXMiniProblems.ProblemPhyXMini0081.answerWavelengthPicometers}
  \uses{def:physics:phyx-mini-0081:phyxminiproblems-problemphyxmini0081-answerchoice}
  The wavelength printed beside each answer label, in picometers
\end{definition}
\begin{proof}
  This is the declared model definition of answer Wavelength Picometers.
\end{proof}
\begin{definition}[Is Unique Closest Displayed Answer]
  \label{def:physics:phyx-mini-0081:phyxminiproblems-problemphyxmini0081-isuniqueclosestdisplayedanswer}
  \lean{PhyXMiniProblems.ProblemPhyXMini0081.IsUniqueClosestDisplayedAnswer}
  \uses{def:physics:phyx-mini-0081:phyxminiproblems-problemphyxmini0081-picometersvalue, def:physics:phyx-mini-0081:phyxminiproblems-problemphyxmini0081-twowavelengthxraybeam, def:physics:phyx-mini-0081:phyxminiproblems-problemphyxmini0081-answerchoice, def:physics:phyx-mini-0081:phyxminiproblems-problemphyxmini0081-answerwavelengthpicometers}
  This def records the Is Unique Closest Displayed Answer component of the problem-specific physical model
\end{definition}
\begin{proof}
  This is the declared model definition of Is Unique Closest Displayed Answer.
\end{proof}
\begin{lemma}[longer Wavelength Picometers bounds]
  \label{lem:physics:phyx-mini-0081:phyxminiproblems-problemphyxmini0081-longerwavelengthpicometers-bounds}
  \lean{PhyXMiniProblems.ProblemPhyXMini0081.longerWavelengthPicometers_bounds}
  \uses{def:physics:phyx-mini-0081:phyxminiproblems-problemphyxmini0081-picometersvalue, def:physics:phyx-mini-0081:phyxminiproblems-problemphyxmini0081-twowavelengthxraybeam, def:physics:phyx-mini-0081:phyxminiproblems-problemphyxmini0081-reflectingcrystal, def:physics:phyx-mini-0081:phyxminiproblems-problemphyxmini0081-angulardiffractionfigure, def:physics:phyx-mini-0081:phyxminiproblems-problemphyxmini0081-matchesproblemandfigurereadouts, def:physics:phyx-mini-0081:phyxminiproblems-problemphyxmini0081-hasphysicaltwolineparameters, def:physics:phyx-mini-0081:phyxminiproblems-problemphyxmini0081-satisfiestwolinebragglaws}
  This lemma records the longer Wavelength Picometers bounds component of the problem-specific physical model
\end{lemma}
\begin{proof}
  The conclusion follows from the governing relations and figure readouts named in the statement of longer Wavelength Picometers bounds.
\end{proof}
% --- Archon declaration topology end ---
% --- Archon physics formalization source end ---
